\subsection{Forced pointed associators and the multiplicity-four frontier}
\label{sec:c4-forced-associator}

The next obstruction uses the associator of a pointed subcategory to restrict
central twists. It is a tensor criterion and applies to an entire family.

\begin{proposition}\label{prop:c4-forced-associator}
Let $m\geq0$, and suppose a fusion category has invertible objects
$1,t,t^2,t^3$, with $t$ of order four, and two self-dual noninvertible
objects $X,Y$ satisfying
\[
 tX=Y=Xt,\qquad tY=X=Yt,\qquad
 X^2=Y^2=1+t^2+m(X+Y),\qquad XY=t+t^3+m(X+Y).
\]
Then the pointed associator has class $2$ in
$H^3(C_4,\mathbb C^\times)\cong\mathbb Z/4\mathbb Z$.
In a pseudo-unitary realization, every simple central summand of
$I(t)$ or $I(t^3)$ has twist satisfying $\theta^4=-1$.
\end{proposition}
\begin{proof}
Put $H=\langle t^2\rangle$. Because $t^2$ fixes a simple object,
the associator restricts trivially to $H$: the associativity scalars
of the action on that object give a two-cochain trivializing the
restriction. Thus its class in $H^3(C_4,\mathbb C^\times)$ is either
$0$ or $2$.

Suppose first that the class is zero. The regular pointed algebra
$A=1\oplus t\oplus t^2\oplus t^3$ exists. Its right-module category
has three simple objects. The invertible orbit contributes $A$;
the orbit $\{X,Y\}$ contributes two simples $M_+,M_-$, each with
underlying object $X\oplus Y$. Indeed its stabilizer is $C_2$,
whose projective irreducible representations are the two ordinary
one-dimensional representations, because $H^2(C_2,\mathbb C^\times)=0$.
This is the usual orbit--stabilizer description of modules over a
regular pointed algebra (equivalently, the classification of modules
over its action groupoid algebra).

Let $d=\operatorname{FPdim}(X)=\operatorname{FPdim}(Y)$, so
$d^2=2+2md$ and $d=m+\sqrt{m^2+2}$. Normalize module dimensions
by $\operatorname{FPdim}(A)=4$; their vector is $(1,d/2,d/2)$.
Every invertible fixes the simple module $A$ and permutes $M_+,M_-$.
On their difference line its action is a character
$\chi:C_4\to\{\pm1\}$, hence $\chi(t^2)=1$.
The action matrix of $X$ is symmetric, since $X$ is self-dual.
Its first row and column are $(0,1,1)$ by free--forgetful adjunction.
The dimension eigenvector therefore forces its remaining block to be
\[
 \begin{pmatrix}a&2m-a\\2m-a&a\end{pmatrix}
 \quad (a\in\mathbb Z).
\]
Thus its eigenvalue on the difference line is the integer
$\lambda=2a-2m$. Applying the fusion relation for $X^2$ on this
line gives
\[
 \lambda^2=2+m(1+\chi(t))\lambda.
\]
For $\chi(t)=-1$ this says $\lambda^2=2$; for $\chi(t)=1$
it says $(\lambda-m)^2=m^2+2$. Both are impossible for integers:
$m^2+2$ lies strictly between consecutive squares for $m\geq1$,
and equals $2$ for $m=0$. Hence class $0$ is impossible.

For completeness use representatives $a,b,c\in\{0,1,2,3\}$ and
standard cyclic cocycles
\[
 \omega_k(a,b,c)=\exp\left(\frac{2\pi i k}{4}
          a\left\lfloor\frac{b+c}{4}\right\rfloor\right).
\]
Their restrictions have $\omega_k(2,2,2)=(-1)^k$, confirming
that the restriction kernel is $\{0,2\}$. The canonical
pseudo-unitary indicator formula gives
$\nu_4(t)=i^k$ and $\nu_4(t^3)=i^{9k}$; for $k=2$ both equal
$-1$. Reversing the cocycle convention inverts these values and
has no effect on this conclusion.

The center indicator formula reads
\[
 \operatorname{FPdim}(\mathcal C)\nu_4(t)
   =\sum_{Z\subset I(t)}[Z:I(t)]\operatorname{FPdim}(Z)\theta_Z^4.
\]
The positive coefficients on the right sum to
$\operatorname{FPdim}(I(t))=\operatorname{FPdim}(\mathcal C)$.
Equality in the triangle inequality forces every $\theta_Z^4=-1$.
The same argument applies to $t^3$; compare
\cite[Proposition~3.8]{SchTwo}.
\end{proof}

\begin{lemma}\label{lem:c4-uniform-module}
Under the hypotheses of Proposition~\ref{prop:c4-forced-associator},
$m$ is even. For the algebra $1\oplus t^2$, its anti-invariant
module action on the two fixed-object coordinates is, up to sign gauge,
\[
 T_t=\begin{pmatrix}0&-1\\1&0\end{pmatrix},\qquad
 T_{t^2}=-I,\qquad T_X=T_Y=0.
\]
In particular both noninvertibles have the anticommuting stabilizer
used in Proposition~\ref{prop:stabilizer-rotation}.
\end{lemma}
\begin{proof}
The involution-module construction gives
$T_t=\left(\begin{smallmatrix}0&\epsilon\\\delta&0\end{smallmatrix}\right)$,
where $\epsilon,\delta\in\{\pm1\}$, and
$T_{t^2}=\epsilon\delta I$. Since $X,Y$ are self-dual,
$T_X$ and $T_Y$ are symmetric integer matrices; they commute with
$T_t$ and satisfy $T_Y=T_tT_X$.
If $\epsilon\delta=1$, commutation makes the diagonal entries
of $T_X$ equal. Its eigenvalue $\lambda$ on the $+1$ eigenspace
of $T_t$ is therefore an integer. But the relation for $X^2$
gives $\lambda^2=2+2m\lambda$, which is impossible as above.
Consequently $\epsilon\delta=-1$. A symmetric matrix commuting
with this skew-symmetric $T_t$ is scalar, so $T_X=aI$.
The symmetry of $T_Y=aT_t$ forces $a=0$. Finally every entry
of $T_X$ is congruent modulo two to the corresponding coefficient
of multiplication by $X$ on $\{X,Y\}$, namely $m$. Hence $m$
is even. Changing the sign of one fixed-object coordinate gives
the displayed sign gauge.
\end{proof}

\begin{theorem}\label{thm:c4-multiplicity-four}
The fusion ring $\RR{4}{6}$ has no categorification over $\mathbb C$.
\end{theorem}
\begin{proof}
Its full tensor has the family in
Proposition~\ref{prop:c4-forced-associator} with $m=4$ and
$(t,t^2,t^3,X,Y)=(b_1,b_5,b_2,b_3,b_4)$.
All simple Frobenius--Perron dimensions belong to
$\{1,4+3\sqrt2\}$. The two-dimension Galois-pseudo-unitarity
theorem therefore reduces general existence to pseudo-unitary
existence, without a pivotality assumption on the original category.

Lemma~\ref{lem:c4-uniform-module} gives exactly one anti-invariant
module representation for the split involution algebra $1\oplus t^2$,
after sign gauge. The finite module checker independently reconstructs it.
The action of $t^2$ is $-I$ on both fixed-object coordinates.
The cubic spaces of $X$ and $Y$ each have dimension four;
Proposition~\ref{prop:stabilizer-rotation} therefore forces
\[
 -2\leq\operatorname{Re}\nu_3(X),
          \operatorname{Re}\nu_3(Y)\leq1.
\]
The exact central-induction calculation has ten low frames and
four second-indicator assignments per frame. All low summands
have twist one. Every high restriction column containing $t$ or
$t^3$ is restricted by Proposition~\ref{prop:c4-forced-associator}
to $\theta^4=-1$. In the normalized relative-trace domain over
$\mathbb Q(\sqrt2)$, this means its fourth coordinate is exactly
$(-1,0)$. Keeping only those trace types is necessary; actual roots
with $\theta^4=-1$ have precisely that coordinate under every
relative Galois embedding. The order-four trace system includes
all root orders with the allowed prime support $\{2,3\}$,
including arbitrarily large orders through their zero trace type.

Each of the resulting forty rational nonnegative systems has a
recorded integral separating vector. Verification reconstructs
the low frames, the unique module branch, the entire joint trace
domain, the four second-indicator assignments, and the forced-twist
variable deletion. For every system $Az=b$, its vector $y$
satisfies $A^Ty\geq0$ and $b^Ty<0$. Therefore no feasible central
induction data, and hence no pseudo-unitary categorification,
exists. This contradicts the general-to-pseudo-unitary reduction.
\end{proof}

\paragraph{Certificate and implementation.}
The tensor-based implementation is
\texttt{software/round3\_smallest.py}. Its verifier never calls
an optimizer; the discovery routine merely proposes exact rays.
The certificate, proof audit, all-census applicability scan,
positive controls and optimized replay are retained in
\texttt{research/round3/smallest/}.
